FedDecay demonstrates sensitivity to the hyper-parameter $\beta$ choice; see Table~\ref{tab: params--sensitivity}.
Specifically, in the case of the FEMNIST data set, runs with $\beta = 0.8$ were prematurely halted due to hyperband stopping, indicating inferior performance compared to smaller $\beta$ choices.
Conversely, the remaining $\beta$ values improved accuracy for new and existing users.
Notably, the FEMNIST data set, comprised of handwritten characters contributed by different authors, showcases inherent data heterogeneity.
This real-world example underscores the substantial user similarities owing to the shared language in which all characters are written.
\begin{table*}[htbp]
    \centering
    \adjustbox{max width=\textwidth}{
        \begin{tabular}{| lc | cc | cc | cc |}
        \hline
             &  & \multicolumn{2}{c|}{FEMNIST (image)} &
                \multicolumn{2}{c|}{SST2 (text)} &
                \multicolumn{2}{c|}{PUBMED (graph)} \\
            Method & $\beta$ & 
                Existing $\Bar{Acc}$ & New $\Bar{Acc}$ &
                Existing $\Bar{Acc}$ & New $\Bar{Acc}$ &
                Existing $\Bar{Acc}$ & New $\Bar{Acc}$ \\
            \hline
            FedSGD & 0.0 & 
                88.51 & 90.55 & 
                73.60 & 80.30 &
                86.70 & 79.14 \\
            FedDecay & 0.2 & 
                \textbf{89.86} & \textbf{91.52} & 
                67.01 & 59.02 & 
                85.67 & 79.68 \\
            FedDecay & 0.4 & 
                89.74 & 90.77 & 
                69.42 & 67.26 & 
                \textbf{87.22} & \textbf{80.39} \\
            FedDecay & 0.6 &
                89.76 & 90.93 & 
                \textbf{78.15} & \textbf{81.01} & 
                86.65 & 79.86 \\
            FedDecay & 0.8 &
                \multicolumn{2}{c|}{Hyperband Stopped} & 
                74.64 & 76.14 & 
                86.59 & 79.32 \\
            FedAvg & 1.0 &
                87.91 & 89.78 & 
                76.54 & 76.80 & 
                86.29 & 79.50 \\
            \hline
        \end{tabular}
    }
    \caption{
        Performance of FedDecay under misspecified values for decay coefficient, $\beta$. 
    }
    \label{tab: params--sensitivity}
\end{table*}


This trend aligns with our assertion in Section~\ref{sec: decay-method-decay}, where we indicated that prioritizing AvgGrad terms (via smaller $\beta$ values) would be beneficial in cases where users share similarities.
Consequently, we observe FedSGD outperforming FedAvg in this context.
However, FedSGD, which does not emphasize AvgGradInner (generalization, rapid personalization), is beaten by FedDecay.
Here, we observe that having a more flexible balance on initial model success, generalization, and fast personalization results in better performance for various choices of $\beta$.

Unlike FEMNIST, heterogeneous users are created from SST2 by partitioning data into 50 clients using Dirichlet allocation with $\alpha = 0.4$.
Many extreme non-iid data sets exist, but $\alpha = 0.4$ results in reasonably non-iid users.
Revisiting the claim discussed in our Section~\ref{sec: decay-method-decay}, we anticipate that an increased emphasis (achieved via larger $\beta$ values) on AvgGradInner is required as non-iid user diversity grows.
In general, larger values of $\beta$ perform better on SST2.
While FedSGD demonstrates superior test-set performance compared to FedAvg, its validation-set accuracy lags behind FedAvg's, preventing it from securing the best run on SST2 for FedAvg in Section~\ref{sec: decay-exp-gen}. 
More importantly, FedDecay balances AvgGrad and AvgGradInner terms, yielding significantly improved metrics over FedAvg.
Notably, FedDecay still has the best overall performance among the three methods.
Lastly, in the case of the PUBMED data set, our method's performance seems resilient to the choice of $\beta$.
However, a straightforward grid search enhances performance for new and existing users over FedAvg and FedSGD.